\section{Contoh Program}
\label{sec:contoh-bab15}

Bagian ini menyajikan dua program yang \emph{sengaja dibuat salah}
beserta versi perbaikannya. Keduanya dirakit tanpa kesalahan sintaks,
sehingga assembler tidak memberi petunjuk apa pun: kesalahan baru
tampak ketika program dijalankan dan ditelusuri.

\subsection{Contoh 1: Perulangan Tak Berujung}

Program pada Kode~\ref{lst:bug-loop} bermaksud menjumlahkan bilangan
1 sampai 5, yang seharusnya menghasilkan \hex{000F} atau \dec{15}.
Kenyataannya program tidak pernah selesai, dan karena baris pencetakan
berada setelah perulangan, program tampak ``menggantung'' tanpa
penjelasan apa pun bagi pemakainya.

\lstinputlisting[
    language=MASM,
    caption={Program ber-bug \file{bug-loop.asm}: perulangan tak
    berujung.},
    label={lst:bug-loop}
]{\CodeDir/bug-loop.asm}

Versi perbaikan pada Kode~\ref{lst:bug-loop-fix} menghapus
\asm{INC CX} dan menambahkan rutin pencetakan dua digit desimal agar
hasilnya dapat dilihat di layar.

\lstinputlisting[
    language=MASM,
    caption={Versi perbaikan \file{bug-loop-fix.asm}.},
    label={lst:bug-loop-fix}
]{\CodeDir/bug-loop-fix.asm}

\begin{session}[caption={Keluaran program
\file{bug-loop-fix.exe}.}, label={lst:keluaran-loop-fix}]
Jumlah 1..5 = 15
\end{session}

\begin{infobox}
Perhatikan bahwa \asm{LOOP} sudah melakukan dua pekerjaan sekaligus:
mengurangi \reg{CX} sebesar satu, lalu melompat bila \reg{CX} belum
nol. Karena itu, penghitung perulangan \emph{tidak perlu} dan
\emph{tidak boleh} diubah sendiri di dalam badan perulangan.
\end{infobox}

\subsection{Contoh 2: Stack yang Tidak Seimbang}

Program pada Kode~\ref{lst:bug-stack} memanggil prosedur
\asm{HITUNG}, lalu bermaksud mencetak pesan penutup. Program berhasil
dirakit, tetapi ketika dijalankan kendali melompat ke alamat yang
tidak dikenal dan pesan penutup tidak pernah muncul. Prosedur
tersebut menekan tiga register (\reg{AX}, \reg{BX}, \reg{CX}) tetapi
hanya memulihkan dua, sehingga \asm{RET} tidak menemukan alamat
kembali pada puncak \eng{stack}.

\lstinputlisting[
    language=MASM,
    caption={Program ber-bug \file{bug-stack.asm}: \eng{stack} tidak
    seimbang.},
    label={lst:bug-stack}
]{\CodeDir/bug-stack.asm}

Versi perbaikan pada Kode~\ref{lst:bug-stack-fix} menambahkan satu
instruksi \asm{POP AX} sehingga jumlah \asm{PUSH} dan \asm{POP}
kembali seimbang.

\lstinputlisting[
    language=MASM,
    caption={Versi perbaikan \file{bug-stack-fix.asm}.},
    label={lst:bug-stack-fix}
]{\CodeDir/bug-stack-fix.asm}

\begin{session}[caption={Keluaran program
\file{bug-stack-fix.exe}.}, label={lst:keluaran-stack-fix}]
Selesai.
\end{session}

\begin{peringatanbox}
Kaidah yang harus dihafal: \emph{setiap} \asm{PUSH} memerlukan
\asm{POP} pasangannya pada jalur eksekusi yang sama. Pada prosedur
yang memakai beberapa jalur keluar, pastikan seluruh jalur
menyeimbangkan \eng{stack} sebelum \asm{RET}. Bila tidak, \asm{RET}
akan melompat ke alamat yang salah dan gejalanya muncul jauh dari
penyebabnya.
\end{peringatanbox}
